\begin{figure}[htbp]
    \centering
    \begin{tikzpicture}[
    font=\sffamily\scriptsize,
    >=Stealth,
    baseBox/.style={
        rectangle,
        rounded corners=2.5pt,
        draw=black!85,
        line width=0.8pt,
        align=center,
        inner sep=4pt
    },
    dados/.style={baseBox, fill=gray!12, minimum width=3.5cm, minimum height=0.6cm},
    personas/.style={baseBox, fill=gray!22, minimum width=7.3cm, minimum height=0.6cm},
    scenario/.style={baseBox, fill=gray!8, minimum width=3.5cm, minimum height=0.6cm},
    llm/.style={baseBox, fill=gray!18, minimum width=7.3cm, minimum height=0.6cm},
    metrics/.style={baseBox, fill=gray!12, minimum width=7.3cm, minimum height=0.6cm},
    seta/.style={->, line width=0.7pt, draw=black!85}
]

    % --- 1. Data sources ---
    \node[dados] (ibge) {Sociodemographic data of\\the Brazilian population};
    \node[dados, right=0.3cm of ibge] (cesop) {Survey of Brazilians'\\perceptions about democracy};

    % --- 2. Synthetic personas ---
    \node[personas, below=0.35cm of $(ibge.south)!0.5!(cesop.south)$, align=center] (personas) {
        2,000 synthetic personas\\(13 sociodemographic attributes)
    };

    % --- 3. Prompt scenarios ---
    \node[scenario, below=0.35cm of personas.south west, anchor=north west, align=center] (c1) {
        C1: attribute selection\\(Essential demography)
    };
    \node[scenario, right=0.3cm of c1, align=center] (c2) {
        C2: natural prose\\(Biographical narrative)
    };
    \node[scenario, below=0.35cm of c1, align=center] (c3) {
        C3: system/user + CoT\\(Structure + reasoning)
    };
    \node[scenario, right=0.3cm of c3, align=center] (c4) {
        C4: few-shot\\(3 personas for calibration)
    };

    % --- 4. Inference and metrics ---
    \node[llm, below=0.4cm of $(c3.south)!0.5!(c4.south)$, align=center] (llms) {
        LLM inference
    };
    \node[metrics, below=0.35cm of llms.south, align=center] (metrics) {
        Evaluation metrics
    };

    % --- Connections ---
    \draw[seta] (ibge.south) -- (personas.north -| ibge.south);
    \draw[seta] (cesop.south) -- (personas.north -| cesop.south);

    \draw[seta] (personas.south -| c1.north) -- (c1.north);
    \draw[seta] (personas.south -| c2.north) -- (c2.north);

    \draw[seta] (c1.south) -- (c3.north);
    \draw[seta] (c2.south) -- (c4.north);

    \draw[seta] (c3.south) |- ++(0, -0.15) -| (llms.north);
    \draw[seta] (c4.south) |- ++(0, -0.15) -| (llms.north);

    \draw[seta] (llms.south) -- (metrics.north);
\end{tikzpicture}

    \caption{Methodological pipeline diagram, illustrating the experimental
        progression across prompt engineering stages.}
    \label{fig:pipeline}
\end{figure}

We structured the experimental investigation into two sequential and
complementary phases. Phase~1 evaluated the cumulative impact of four
prompt engineering strategies on distributional convergence, using a single
high-cardinality anchor question (P02) to isolate prompting effects.
Building upon the optimal prompt configuration identified in the first stage,
Phase~2 conducted an architectural benchmark comparing Brazilian
Portuguese-centric models against general-purpose global baselines across three
distinct thematic dimensions (P02, P03, and P04) to assess sociocultural fidelity
and generalization. The subsequent subsections detail the empirical reference
data, the synthetic population generation, the evaluated architectures, and
the experimental prompt scenarios.

\subsection{Empirical Reference Survey: Perceptions of Democracy}
\label{subsec:database}

We adopted the public opinion survey ``Perceptions of Brazilians Regarding
Democracy''~\cite{cesop04832} as our empirical ground truth. Fielded in September
2023 by the Center for Public Opinion Research (CESOP/UNICAMP) in partnership
with IPEC, the survey sampled 2,000 voting-age citizens across 127 Brazilian
municipalities using a three-stage probabilistic cluster design with quota
allocation based on official census indicators (margin of error of $\pm 2\%$ at a
95\% confidence level).

From the survey questionnaire, we selected three evaluative attitudinal
questions for comparative evaluation: P02 (Policy Priorities: ``Which of
these proposals do you think should be the priority of a politician?'',
featuring 12 substantive response alternatives); P03 (Fake News
Countermeasures: ``Which of these options do you believe could contribute to
combating the spread of fake news?'', with 6 substantive alternatives on regulatory and
institutional interventions); and P04 (Political Participation: ``Would you
say that you have a strong desire, some desire, or no desire to participate in
the political life of your city?'', offering 3 substantive alternatives measuring local civic
engagement).

Our selection includes all evaluative attitudinal items from the instrument,
deliberately omitting item P01 due to its purely retrospective nature as an
electoral recall check on 2022 voting behavior. This tripartite set provides
thematic breadth, spanning material policy agendas, institutional responses
to misinformation, and subjective civic disposition, while also varying in
response structure. By varying response cardinality across $12$, $6$, and
$3$ substantive categories, these questions allow us to evaluate JSD convergence under
different levels of response entropy and potential social desirability effects.

\subsection{Generation of the Synthetic Population Grounded in IBGE}
\label{subsec:silicon_profiles}

We generated 2,000 synthetic individual profiles using stratified
probabilistic sampling based on official distributions from the IBGE (Brazilian
Institute of Geography and Statistics), retrieved via the SIDRA
API~\cite{ibge_sidra_api}. Each persona is characterized by up to 13 attributes
spanning demographic, economic, and social dimensions, parameterized from the
2022 Demographic Census~\cite{ibge_censo_2022} (gender, age bracket, race,
region, zone, and religion), Continuous PNAD 2024~\cite{ibge_pnadc} (education,
employment, and \textit{per capita} income), the National Health Survey (PNS
2019)~\cite{ibge_pns_2019}, Civil Registry~\cite{ibge_registro_civil}, and
IPCA inflation indicators~\cite{ibge_ipca}.

Grounding the synthetic population directly in official IBGE microdata
guarantees reproducible probabilistic sampling across Brazilian federative units
and socioeconomic strata, ensuring the methodology remains transferable to
scenarios where prior survey ground truth does not pre-exist.

\subsection{Evaluated Large Language Model Architectures}
\label{subsec:language_models}

We compared four LLM architectures throughout our experiments, as summarized in
Table~\ref{tab:modelos}.

\begin{table}[htbp]
    \centering
    \caption{Large Language Models evaluated in the study.}
    \label{tab:modelos}
    \begin{tabular}{@{}llll@{}}
        \toprule
        \textbf{Model} & \textbf{Provider} & \textbf{Type} & \textbf{Origin} \\
        \midrule
        Sabiá-4 & Maritaca AI & API & Brazilian \\
        Sabiazinho-4 & Maritaca AI & API & Brazilian \\
        GPT-5-mini & OpenAI & API & Global \\
        LLaMA~3.2~3B & Meta & Local & Global \\
        \bottomrule
    \end{tabular}
\end{table}

The choice of Brazilian models (Sabiá-4 and Sabiazinho-4)
was grounded in the hypothesis that models pre-trained with an emphasis on the
Portuguese language and Brazilian sociocultural context could reproduce local
public opinion with higher fidelity. We selected GPT-5-mini as a
global baseline in the equivalent cost-performance tier, consistent with
the technical benchmark released by Maritaca AI for the Sabiá
architecture~\cite{abonizio_sabia-3_2024}.

\subsection{Experimental Prompt Engineering Strategies}
\label{subsec:cenarios}

We executed a cumulative progression of four experimental scenarios applying
incremental refinements to the prompt design strategy. In all scenarios, we
utilized 2,000 synthetic personas and initially focused on question
P02 (Policy Priorities). To mitigate the impact of probabilistic
sampling during inference, we repeated each prompt execution five times and
adopted the majority vote outcome for evaluation.

\subsubsection{Scenario 1: Feature Selection}
\label{subsubsec:cenario1}

In the first scenario, we applied a feature selection technique by using the
LLM itself to identify the most predictive variables among the
13 original attributes for the question at hand. For P02,
the model selected seven attributes: education,
\textit{per capita} income, region, age bracket,
self-rated health, gender, and marital status.

We structured the prompt as a single user message, presenting the profile as a
list within XML tags and requesting solely the letter corresponding to the
chosen alternative:

\begin{quote}
    \small
    \textit{``You must think and respond as a Brazilian citizen with the profile
    below. [\ldots] Consider the socioeconomic and demographic conditions of the
    profile to form your opinion on the question and respond with the letter of
    the alternative that best represents your opinion. [\ldots] My answer is''}
\end{quote}

We set the temperature to 0.7 and \texttt{max\_tokens} to
100, enforcing structured schema output (\texttt{response\_format}) to
ensure automated extraction.

\subsubsection{Scenario 2: Natural Prose Biographical Representation}
\label{subsubsec:cenario2}

As a natural evolution of the previous scenario, we utilized role-playing
techniques to assess how the model interpreted the presented profile. We
rephrased the prompt from a list of attributes into a natural language
descriptive paragraph:

\begin{quote}
    \small
    \textit{``Act as a Brazilian individual of gender [gender], from the [region]
    region, aged between [age bracket], [marital status]. Your education level is
    [education] and your individual income is [income]. Your self-rated health is
    [health].''}
\end{quote}

The remaining parameters (temperature, response format, and model) were kept
identical to Scenario~1.

\subsubsection{Scenario 3: Structural Role Separation and Chain of Thought (CoT)}
\label{subsubsec:cenario3}

In this scenario, we introduced two complementary structural modifications.
First, for \textit{role separation}, we allocated the persona profile and
behavioral guidelines exclusively to the \textit{system prompt}, providing the
question and alternatives in the \textit{user prompt}, reinforced by explicit
instructions to prevent breaking character: \textit{``Your task is to answer
public opinion questions exactly as this person would answer --- not as you
believe is correct, nor as an expert would answer.''} Second, to implement
\textit{Chain of Thought (CoT)}, we instructed the model to reflect briefly on how
an individual with the described profile would experience the issues in daily life
before selecting the final alternative, enforcing a structured response schema
with two fields: \texttt{reasoning} and \texttt{answer}.

We set the temperature to $0.9$ and adjusted \texttt{max\_tokens} to 200 to
allow for intermediate reasoning generation. Additionally, we reintroduced all
13 original profile attributes into the \textit{system prompt}.

\subsubsection{Scenario 4: Contextual Few-Shot Calibration}
\label{subsubsec:cenario4}

In the final scenario, we maintained the system/user prompt structure from
Scenario~3 and added three contrasting calibration personas as prior conversational
turns (assistant/user) preceding the target question. Specifically, the calibration
exemplars comprised: (1)~a young, low-income Black woman from the urban Northeast
choosing \textit{``Combat prejudice''} (alternative~b); (2)~an older, rural White
Evangelical man from the South choosing \textit{``Preserve family-oriented values''}
(alternative~g); and (3)~a middle-aged, college-educated divorced woman from the
urban Southeast choosing \textit{``Reduce social inequalities''} (alternative~a).

We constructed these three examples from qualitative analysis of
justifications generated by the LLM in previous runs (Scenario~3),
reusing them across all 2,000 simulations. We configured the structured
output to require three fields: \texttt{top\_concern} (primary daily concern),
\texttt{why\_this\_answer} (perspective justification), and \texttt{answer}
(alternative letter).

We appended an instruction to the end of the system prompt to mitigate inherent
training bias:

\begin{quote}
    \small
    \textit{``ATTENTION: In the real Brazil, surveys show that different
    Brazilians choose very different priorities. Prejudice, Inequality, Family
    values, Education, and Political participation are equally legitimate and
    frequent priorities. Health and Jobs are important to everyone, but they are
    not automatically the answer for every persona.''}
\end{quote}

To evaluate distributional convergence between simulated distributions ($P$)
and empirical CESOP targets ($Q$), response category counts were normalized to
unit sum ($\sum p_i = \sum q_i = 1$) and evaluated using base-2 Jensen-Shannon
Divergence (JSD)~\cite{lin_divergence_1991}. Bounded within $[0, 1]$ bit,
$\text{JSD} = 0$ denotes identical distributions. In line with computational
social simulation standards~\cite{miranda_simulating_2025}, we adopt $\text{JSD}
< 0.15$ as the threshold indicating close distributional convergence, whereas
values exceeding $0.50$ indicate severe empirical divergence.

